\chapter{Application Hugging Face Spaces}

\section{Objectif}
L'application fournit une interface utilisateur pour generer des narratives analogues
de maniere controlee.

\section{Fonctionnement}
L'utilisateur selectionne:
\begin{itemize}[leftmargin=2em]
  \item une entite source,
  \item une langue source,
  \item une instruction d'adaptation,
  \item un seuil de similarite,
  \item un modele de similarite (allMiniLM FT, MPNET FT, ou custom).
\end{itemize}

Pour chaque tentative:
\begin{enumerate}[leftmargin=2em]
  \item le LLM genere un candidat narratif,
  \item le sentence-encoder calcule la similarite avec la source,
  \item si \(score < seuil\), le systeme relance une tentative en tenant compte
  de l'historique des echecs.
\end{enumerate}

\section{Sorties affichees dans l'UI}
\begin{itemize}[leftmargin=2em]
  \item \textbf{Run Summary}: accepted, best score, nombre de tentatives, tentative acceptee;
  \item \textbf{Best Candidate Narrative}: meilleur narratif retenu;
  \item \textbf{Attempts}: tableau minimal avec \texttt{attempt}, \texttt{score}, \texttt{accepted};
  \item \textbf{Full Result JSON}: sortie complete pour audit/reproductibilite;
  \item \textbf{Execution Logs}: traces d'execution.
\end{itemize}

